% Extracto íntegro por residencia del propietario científico definitivo de masas.
\section{Internal domain and realization problem}

The General Mass Law does not start from a table of names or a list of
experimental values. Its internal domain is the finite chain
\[
 \mathcal R_{324}\longrightarrow \mathcal C_{81}
 \longrightarrow \mathcal F_{56}\longrightarrow \mathcal M_{13},
\]
where \(\mathcal R_{324}\) contains the \(4\) orientations of each of the
\(81\) cells, \(\mathcal F_{56}\) groups equivalent route histories, and
\(\mathcal M_{13}\) preserves the operator multisections. The null sector
bifurcates before the positive character is applied. Therefore, neither an
abbreviated historical table nor a taxonomy of physical names can replace this
domain.

Let \(\mathcal P\) be the set of individualized physical states. An elementary
state is represented by the descriptor
\[
 X=(\mathfrak c,\rho,g,f,q,\epsilon,\varsigma,\mathfrak a),
\]
which preserves representation class \(\mathfrak c\), internal representation
\(\rho\), generation \(g\), flavor \(f\), charge \(q\), sheet orientation
\(\epsilon\), statistical sector \(\varsigma\), and coupling data
\(\mathfrak a\). A composite state adds ordered constituents, composition
tree, common boundary, and collective quantum numbers. A resonance
adds the open channels and the causal prescription for the resolvent.

The physical realization must address three distinct questions:

\begin{enumerate}
\item which fibre, orbit, or multisection of the atlas corresponds to the descriptor;
\item which positive operator publishes the mass on that fiber;
\item which coupling produces a scalar, a multiplet, a distribution, or
      a complex pole.
\end{enumerate}

Confusing these questions introduces a retrospective selector. The correct
answer is obtained through a natural relation, a fiber operator and a
spectral compression, in that order.

\section{Fiber product between descriptors and routes}

Let \(\mathcal D\) be the finite space of internal data already constructed before the
contrast: representation, orientation, memory, boundary, charge, flavor,
color, generation, and statistical sector. Denote by
\[
 d_{\mathcal P}:\mathcal P\longrightarrow\mathcal D,
 \qquad
 d_{\mathcal R}:\mathcal R_{324}\longrightarrow\mathcal D
\]
the physical descriptor maps and the genealogical descriptor. The object
of realization is defined by the fiber product
\[
 \mathcal J
 =\mathcal P\times_{\mathcal D}\mathcal R_{324}
 =\{(X,\gamma):d_{\mathcal P}(X)=d_{\mathcal R}(\gamma)\}.
\]
The fiber of a species is
\[
\mathscr S(X)=\{\gamma\in\mathcal R_{324}:(X,\gamma)\in\mathcal J\}.
\]
A descriptor (X) is called realizable by the atlas when
\[
 d_{\mathcal P}(X)\in\operatorname{im}d_{\mathcal R}.
\]
This condition is equivalent to \(\mathscr S(X)\ne\varnothing\). It does not follow from merely writing the fiber product: it constitutes the exact existence condition that must be verified for each physical descriptor. The value of \(\mathscr S\) may be a route, an orbit, or a multisection. A point section is not forced when the internal symmetry requires preservation of a multiplicity.

\begin{definition}[Admissible realization relation]
The relation \(\mathscr S\) is admissible when it simultaneously satisfies:
\begin{enumerate}
\item depends solely on data constructed before the external table is opened;
\item commutes with particle--antiparticle conjugation;
\item is equivariant under cellular inversion and the internal
      gauge actions;
\item it is compatible with refinement, nonadic return and dodecaphasic
      reading;
\item preserves every multiplicity not eliminated by a physical coupling;
\item remains invariant under permuting or replacing the external values of
      mass and width.
\end{enumerate}
\end{definition}

\begin{theorem}[Prospective realization by fibered product]
If \(d_{\mathcal P}\) and \(d_{\mathcal R}\) are equivariant with respect to a symmetry \(G\) and do not depend on external magnitudes, then \(\mathcal J\) is a \(G\)-subobject of \(\mathcal P\times\mathcal R_{324}\). For every realizable descriptor \(X\), the nonempty fiber \(\mathscr S(X)\) is \(G_X\)-stable, and the assignment \(X\mapsto\mathscr S(X)\) is independent of the result subsequently compared.
\end{theorem}

\begin{proof}
For \((X,\gamma)\in\mathcal J\) one has
\(d_{\mathcal P}(X)=d_{\mathcal R}(\gamma)\). Equivariance gives
\[
 d_{\mathcal P}(gX)=g\,d_{\mathcal P}(X)
 =g\,d_{\mathcal R}(\gamma)=d_{\mathcal R}(g\gamma),
\]
and, consequently, \((gX,g\gamma)\in\mathcal J\). If \(g\in G_X\), the
second component remains in \(\mathscr S(X)\). Independence of the
external values is immediate because none of them belongs to the domain
of the two maps of the fiber product.
\end{proof}

The electron constitutes the central exception. Cell inversion has a unique fixed point, the cell \((5,5)\), and its physical fiber therefore has rank one. The remaining stable multisections do not, in general, admit an equivariant point selector. This difference is not a deficiency of the atlas: it is the mathematical distinction between a central section and a species realized as an orbit or multisection.

\section{Intrinsic spectral classification of the species}

Let \(\mathcal A_{\rm rut}\) be the finite commutative algebra generated by the
internal route observables: signature \(\mathbb Z^5\), character, chronology,
\(K_D\), \(K_\Omega\), orientation, memory, and compatible readers. In a
fiber \(F\), one defines
\[
 \gamma\sim_{\rm sp}\eta
 \quad\Longleftrightarrow\quad
 a(\gamma)=a(\eta)\quad\text{for every }a\in\mathcal A_{\rm rut}.
\]
For each class \(\lambda\in F/\!\sim_{\rm sp}\),
\[
 P_\lambda=\sum_{\gamma\in\lambda}
 |\gamma\rangle\langle\gamma|
\]
is a spectral projector; the \(P_\lambda\) are orthogonal and
\(\sum_\lambda P_\lambda=I_F\).

\begin{theorem}[Maximum internal individualization]
Every classifier constructed exclusively from observables of \(\mathcal A_{\rm rut}\) factors uniquely through the quotient \(F/\!\sim_{\rm sp}\). Consequently, two descriptors with distinct complete signatures belong to distinguishable sectors; two descriptors that agree for all generators of the algebra cannot be separated by a rule constructed from those same observables.
\end{theorem}

\begin{proof}
The characters of the finite commutative algebra separate exactly its
joint spectrum. A functional classifier of its generators is constant
on each fiber of the joint spectral map and therefore factors through the
quotient. If two routes have the same character on the entire algebra, no
element of it can separate them.
\end{proof}

This theorem closes the logical limit of the classification. When a
physical representation requires two states belonging to the same
spectral class to be distinguished, it must exhibit an additional observable or a coupling that
enlarges \(\mathcal A_{\rm rut}\). The absence of such an observable does not authorize
the external mass to be used as a selecting label.

In the charged leptonic sector, the descriptors distinguish the central fiber,
the depth-\(G_2\) multisection, and the depth-
\(G_4\) multisection. The first produces the unique electronic line; the other two
produce the operators associated with the muonic and tauonic realizations. If
the coupling preserves more than one block, the correct prediction is its
spectrum; the physical label does not authorize choosing one of its eigenvalues by
numerical proximity.


